\documentclass[11pt]{article}
\usepackage{../common/indexus}
\usepackage{../figkit/ixfig}

\begin{document}

\title{Indexus: The Client}
\author{\ixbyline}
\hypersetup{pdftitle={Indexus: The Client}, pdfauthor={Rémi Henriot (AI Co-authored)}}
\date{November 2026}
\maketitle

\begin{abstract}
\noindent \textbf{Abstract. }The Distributed Prefix Index Protocol (DPIP) lets a reader refine a view one prefix at a time~[\ixref{2}]. This paper follows the Indexus reference client. A client turns coordinates into locations, reads the sums of a zone's children in a compact reply, draws a map at any zoom by reading only the zones it lacks, searches around a point with the incremental algorithm of Henriot's 2024 paper~[\ixref{1}], and joins collections of different producers cell by cell. Everything here is implemented in the reference client.
\vspace{1em}\end{abstract}

\section{Introduction}
A reader of the index wants two things. A map wants totals: how many records, and how much, in each part of what it shows, at the size given by its zoom. A feed wants records: the nearest ones first, then the next, as the user scrolls. Both read the same zones and the same sums. The map reads counts and totals from the index. It can also compute a displayed quantity, such as an average, from those totals.

We follow the whitepaper's words: a \emph{record} is one entry, a \emph{collection} a named set of records sharing a space, a \emph{zone} the records whose location starts with the same characters, and the \emph{sum} of a zone its count and the totals of its numbers.

\paragraph{A map as a worked example.}
Suppose the view contains one city cell with 60 houses. At the next zoom the client asks for that cell's children, not for all 60 houses. It receives counts such as 15, 3, 33 and 9, draws the four cells, and reads deeper only where the screen needs detail. Nearby search has a different goal: eventually return individual houses in distance order. Both operations start from the same summaries.

\paragraph{Companion papers.}
\ixsection{whitepaper}{space}{The whitepaper's space example} introduces the geometry.
\ixsection{network}{authority}{The network paper's authority section} explains why
an outdated route cannot grant a second writer. The broader summary operations
and their arithmetic limits are in \ixsection{properties}{algebra}{the properties paper}.

\section{Spaces}\hypertarget{ix:spaces}{}
A space is a list of dimensions. A dimension has a range and a distance: a sphere has latitude, longitude and the distance along the surface; a line has a start, an end and the difference of two values. When the mask allocates bits equally and $n$ divides 6, a space of $n$ coordinates cuts each of them $6/n$ times per character: a map of latitude and longitude cuts each three times per character.

\begin{figure}[H]
    \centering
    \ixfig[12cm]{../figkit/orig/masks.tikz}
    \caption{Two subdivisions of the same space}
\end{figure}

The client computes a location as follows. The mask gives the dimension cut by each bit, and the offset is the collection's name read as numbers.

\begin{algorithm}[H]
\caption{Location of a point}
\begin{algorithmic}[1]
\State $bounds \gets$ the ranges of the space; $location \gets$ ``''
\For{$l = 0$ \textbf{to} $precision - 1$}
  \State $bits \gets 0$
  \For{$n = 0$ \textbf{to} $5$}
    \State $a \gets mask(l, n)$; $mid \gets$ middle of $bounds[a]$
    \If{$point[a] > mid$} $bits \gets 2\,bits + 1$; set $bounds[a]$ to $[mid, max]$
    \Else\ $bits \gets 2\,bits$; set $bounds[a]$ to $[min, mid]$
    \EndIf
  \EndFor
  \State $location \gets location + alphabet[(offset(l) + bits) \bmod 64]$
\EndFor
\end{algorithmic}
\end{algorithm}

Decoding replays the same halvings after removing the offset, and gives back the \emph{cell}: the box of coordinates that the location covers. With base 4 the rule is the same with 2 bits per character: the point at latitude $-60$ and longitude $-45$, with an offset BC, has the location BD and the cell latitude $[-90,-45]$ by longitude $[-90,0]$. With the real alphabet, Paris has a four-character location whose cell is about 4.9 by 6.5 kilometres.

\section{Zones}
A location is a path in the tree. Each prefix names a zone, and removing the last character gives its parent. A zone is identified by its collection and location. The specified routing key starts with the location bits and preserves the remaining bits of the collection name. When a prefix ends inside a byte, implementations must agree on those remaining bits. The \ixpaper{network} discusses the observed discrepancy. The reader retains the full collection and location when requesting the zone.

\begin{figure}[H]
    \centering
    \ixfig[10cm]{../figkit/orig/encoded.tikz}
    \caption{Encoded tree (BASE4)}
\end{figure}

The address labels a delivery destination. Removing the collection offset recovers the underlying cell used later for composition.

\section{Reading a zone}\hypertarget{ix:zone-read}{}
A read of a zone answers the sums of its children: one row for the records placed exactly at the zone, and one row per non-empty child, with its count and the total of each number. It has at most 65 rows. When every record of a child lies in a single place, the row is given under that place, down to four characters below the zone; the client can then refine a map without reading further.

\begin{figure}[H]
    \centering
    \ixnew{figures/sums.tikz}
    \caption{Reading zone @ returns its children A, B, C and D with their counts; reading zone C returns the records placed at C and its children CB, CC and CD.}
\end{figure}

Rows travel as columns of numbers: the keys, the counts, then each metric. A browser reads them where they arrive and hands them to its graphics processor without building one object per row. A client may ask for some metrics only; the sums are still those of whole records.

In this simplified example, the answer for C counts 33 records. A closer view can separate four placed exactly at C from 9 in CB, 10 in CC and 10 in CD. The four rows add back to 33. Omitting the row at C would drop four houses from the map.

\paragraph{Choosing a read.}
The map can still display its four cells during a refresh. Checking a new house after acknowledgement requires a fresh server answer. Section~\ref{sec:freshness} follows that house through both reads.

\section{Finding nodes}
In a gateway-backed application, \texttt{GatewaySession} obtains the network settings and a client credential, then opens a \texttt{Network} connection. Once connected, the map and the nearby feed can share this connection while keeping their own traversal state. A client can let one entry node route its reads or ask the nearest known node directly. A redirect adds a hop if a zone has moved. The client can group up to 128 zones per request.

\section{Drawing a map}\hypertarget{ix:map}{}
A map shows a region at a zoom. The zoom sets the size of the cells to draw. The client keeps the zones it has already read and descends them: a cell outside the view is skipped, a cell fine enough is drawn, a cell too coarse is replaced by its children if they are known. The children it lacks are exactly the zones it reads next.

\begin{figure}[H]
    \centering
    \ixnew{figures/zoom.tikz}
    \caption{Map refinement. The client reads @, then C, then the four children of C.}
\end{figure}

\paragraph{A cell on the screen and a zone on the wire.}
On a two-dimensional grid that alternates the axes, one BASE64 character contains three latitude and three longitude cuts. One child-summary block can therefore support three successive map refinements: the reader groups its rows by the first two, four, then six geometric bits after undoing the offset. Further key digits present in the block can also place rows in finer bins. The worker selects the visible cells and the blocks still needed. WebGPU groups the selected rows for drawing. Zooming one step need not trigger a network read.

Suppose a parent attributes 30 houses to a branch, but the child block currently accounts for 24. Replacing the visible branch with that breakdown would make six houses disappear. The \href{https://github.com/indexus/sdk/blob/6e1bba2/src/walk/index.js}{viewport traversal} can keep the coarser cell in the finer view. Matching counts allow refinement; they cannot establish that both observations describe the same instant. On an initial load there may be no cached parent to show. The map reads a margin around the viewport and prepares neighbouring resolutions so that a small pan or zoom can reuse the work.

\section{Searching around a point}
To find the nearest places, the reader starts at an origin and opens zones from coarse to fine, as in Henriot's 2024 search algorithm~[\ixref{1}]. For each zone $Z$ it computes a lower bound $b(Z)$ on the distance from the origin to any record in that zone. If the nearest unopened zone starts at 5\,km, a record already found at 3\,km can be shown. One at 8\,km must wait because the unopened zone may contain a closer place. Counts help decide how much to fetch. The client reads selected zones in batches and can fetch the records of a small zone directly. The distance bound determines which results can be shown in order.

\paragraph{What the next page remembers.}
The client keeps the search in memory. Each layer retains the zones still to open, its distance bound, and the entries already selected or read. A page returns up to the requested number of terminal records. Any others already fetched stay available. A second call on the same search advances these layers and can return a page without another network request. Distinct records can share both a location and a distance. Each keeps its own identifier.

Consider a search between 5 and 10\,km from the origin. A coarse cell containing the origin has a lower bound of zero, yet it may contain eligible records. The reader opens intersecting cells and tests the actual records against both limits. In the \href{https://github.com/indexus/sdk/blob/6e1bba2/src/local/annular-filter.test.mjs}{one-dimensional pagination test}, four records lie at distances 1, 6, 8 and 12 from the origin. With a page size of one, the first call returns the record at 6, the second returns the one at 8. Across both calls the reader fetches the root and the coarse cell once each.

\paragraph{When to start again.}
Changing the origin, filters or collections changes which records qualify and their order. Start a new search for those inputs. The old traversal can be \href{https://github.com/indexus/sdk/blob/6e1bba2/src/local/cancellation.test.mjs}{cancelled}. Continuation lives in that client instance. Reloading the application loses it. Ordering assumes the required cells can be read from stable data. A zone read can fail and the \href{https://github.com/indexus/sdk/blob/6e1bba2/src/local/exploration.js}{reference traversal} can continue past it. Such a page can omit a nearer record. Writes between pages can also change the ranking, since separate reads have no shared snapshot.

\begin{figure}[H]
    \centering
    \ixfig[14cm]{../figkit/orig/layers.tikz}
    \caption{Layers and concentric circles. The \href{https://indexus.io/papers/animations/search-algorithm.gif}{original 20-frame animation} follows the search through three depths of the tree.}
\end{figure}

\section{Combining collections}\hypertarget{ix:composition}{}
In the 2024 search paper~[\ixref{1}], a mask selects which coordinate to bisect next, and a collection-specific offset changes the letters assigned to the resulting cells. The earlier paper also described searching across several collections in one traversal. Here the operation is arithmetic on cell summaries. The encoding lets two producers describe the same place with different strings. An address is useful to find the data. To compare two cell measurements, the reader instead decodes their integer grid coordinates and resolution.

\begin{figure}[H]
    \centering
    \ixnew{figures/compose.tikz}
    \caption{One shaded cell on two BASE4 grids. Its two low bits are 00. Adding offsets B and D encodes it as B and D respectively. Both decode to cell $(0,0)$ at the same resolution.}
\end{figure}

\paragraph{What has to agree.}
The axes must have the same meaning, order, units and bounds, and the subdivision schedule must give the cells the same footprints. The \href{https://github.com/indexus/sdk/blob/6e1bba2/src/map/heat/compose.js}{map composer} compares the ordered dimension types and arguments. Its ordinary collection constructor uses the same cyclic mask for those axes. The caller still supplies the meaning of each field and its units. With a different coordinate system or subdivision schedule, an application must first express the data on a common grid. Dropping the offsets alone does not reproject it.

\paragraph{One calculation, cell by cell.}
Suppose two producers publish daytime and night-time temperatures for the same date. In one cell, two daytime readings sum to $58\,^{\circ}\mathrm{C}$ and three night-time readings sum to $69\,^{\circ}\mathrm{C}$. Their means are 29 and 23. A map of positive daytime warming displays $\max(0,29-23)=6\,^{\circ}\mathrm{C}$. The reader can evaluate the expression from four numbers, even if the producers sampled different sensors. Matching a particular sensor reading would require a shared identifier and observation time.

The first field sets the cells drawn. For each of those cells the reader looks up the others, up to four fields in all. If the only night-time value is the mean of an enclosing cell, 23 describes that wider area; the displayed difference has correspondingly coarser support. If there are only finer cells, the composer currently leaves that field missing. Without a declared substitute it hides the output cell. Treating missing temperature as zero would invent a 29-degree difference. The \href{https://github.com/indexus/sdk/blob/6e1bba2/src/gpu/combine.test.mjs}{kernel tests} exercise these three cases against a reference calculation. The \ixsection{properties}{anglet}{Anglet example} follows a different use of several collections in a single view.

\section{Keeping data}
The client keeps the answers it has read in memory, within a budget of 256~MB by default. When the budget is exceeded, the least recently used answers leave, except those the current view needs. The same answers are copied once to the graphics processor, which adds the rows of each cell to draw and writes one instance per cell. If the graphics device is lost, it is rebuilt from memory without reading the network again.

\section{Freshness}\label{sec:freshness}\hypertarget{ix:freshness}{}
Suppose a map has cached 60 houses and a producer then receives an acknowledgement for a new house. There are two places where the older count may remain visible. The map may reuse its answer. If it asks the server for an ordinary summary, the parent there may still have only the older child version. Neither delay has a fixed upper bound during disconnection.

\begin{figure}[H]
    \centering
    \ixnew{figures/refresh.tikz}
    \caption{One new house after a cached count of 60. A fresh server read must cover the current holders.}
\end{figure}
For a check after the ACK, send a fresh request to the server with \texttt{refresh=true}. When using the \href{https://github.com/indexus/sdk/blob/83f2db1ad52ac695ab01cb2ca188f23300271bd7/src/network/index.js#L107-L114}{reference SDK}, set \texttt{refreshTtlMs} to zero if this check must bypass the default five-second reuse of a recent result. In normal mode the server must cover the required holders and any receivers of transferred zones, or report failure. With no other changes, a successful fresh read returns 61. If the holder cannot be reached, the client can surface the failure. A reader that explicitly permits \texttt{stale-if-error} can instead receive an older observation within its age limit. That age runs from the node's observation. A batch from \texttt{/sets} promises no common observation time across prefixes or collections. The server's \href{https://github.com/indexus/core/blob/5188b61/docs/protocol/read-modes.md}{read policies} define the distinction. The \ixsection{network}{read-consistency}{network paper} explains coverage during a move and the limits of this guarantee.

\section{Calculations}
\textit{Size of an answer.} With one 32-bit suffix word per key, the count and key take 8 bytes per child. Four default 32-bit metrics add 16 bytes, giving $64\times24=1\,536$ bytes of column data for 64 children. One such root block with its headers is 1\,572 bytes before compression. A row at the parent, a second key depth, redirects or 64-bit metrics change this size.

\textit{Zones read for a view.} With latitude and longitude, each character divides each axis by 8. A view covering a fraction $w$ by $h$ of the world, read at a depth $D$, needs about $\lceil w \cdot 8^{D} \rceil \times \lceil h \cdot 8^{D} \rceil$ zones, fewer where the map is empty.

\textit{Search.} A step reads the zones whose remaining distance bounds still intersect the desired results. A denser frontier can require several batches at the same depth. Both the number of relevant zones and how they are distributed across nodes affect the request count.

\section{Conclusion}
For a map the client draws summaries and refines the visible cells; for a feed it ultimately reads individual records. A map can also combine collections at a common cell even when different producers use different offsets. A reader checking a newly acknowledged write can request a fresh result and handle an unavailable holder explicitly.

\Needspace{5\baselineskip}
\paragraph{Contributions.} The search algorithm and the original drawings are Rémi Henriot's. AI-assisted review traced the SDK cache and helped develop the example of a fresh read after an acknowledged write. The \ixsection{history}{ai}{project history} dates the contributions of the tools used.

\section*{References}
\begin{enumerate}[label={[\arabic*]},leftmargin=2em,itemsep=1pt]
  \item \hypertarget{reference:1}{}R. Henriot. \ixbib{search2024}. 2024.
  \item \hypertarget{reference:2}{}R. Henriot. \ixbib{whitepaper}. 2026.
\end{enumerate}

\end{document}
